% Compile with: xelatex article30k_m2pro_en.tex; bibtex article30k_m2pro_en; xelatex ...; xelatex ...
\documentclass[11pt,a4paper]{article}

\usepackage[margin=2.35cm]{geometry}
\usepackage{fontspec}
\setmainfont{Arial}
\usepackage[main=english]{babel}
\usepackage{microtype}
\usepackage{amsmath}
\usepackage{booktabs}
\usepackage{tabularx}
\usepackage{array}
\usepackage{enumitem}
\usepackage{caption}
\usepackage{hyperref}
\usepackage{xurl}
\usepackage[numbers,sort&compress]{natbib}

\hypersetup{
  colorlinks=true,
  linkcolor=black,
  citecolor=black,
  urlcolor=blue,
  pdftitle={Multi-task Classification of Review Sentiment and Provenance},
  pdfauthor={Author to be specified before submission}
}
\newcolumntype{Y}{>{\raggedright\arraybackslash}X}
\newcommand{\repo}[1]{\nolinkurl{#1}}
\setlength{\emergencystretch}{2em}

\title{Multi-task Classification of Review Sentiment and Provenance with XLM-RoBERTa:\\
A Reproducible Study on a Partially Labelled Corpus of 30,000 Records}
\author{\textbf{[Author name]}\\
\textit{[Programme, university]}\\
\texttt{[author-email]}}
\date{September 2026}

\begin{document}
\maketitle

\begin{abstract}
Online reviews contain both evaluative judgements and signals about how a text was produced, yet these properties are often modelled separately. This study tests whether shared training of one language encoder can address two tasks without sacrificing quality: three-way sentiment classification and discrimination between human-written and AI-generated reviews. A corpus of 30,000 records was assembled from four open datasets; provenance labels are available only for 7,000 hotel reviews from MAiDE-up. TF--IDF with logistic regression, two single-task and one multi-task XLM-RoBERTa systems are compared. Transformer variants were trained with five random seeds, while a fixed split protected against normalized-text overlap across partitions. On the test set, the multi-task model reached Macro-F1 of $0.7709\pm0.0087$ for sentiment and $0.9526\pm0.0094$ for provenance; the corresponding classical-baseline scores were $0.7161$ and $0.8998$. Differences from single-task Transformer systems were inconclusive: $-0.0013$ for sentiment and $-0.0046$ for provenance. Within this protocol, multi-task learning is therefore a compact alternative to two separate models, but it is not evidence of positive transfer over single-task training. The findings apply to a controlled within-protocol setting and do not establish universal fake-review detection.

\medskip
\noindent\textbf{Keywords:} sentiment analysis; multi-task learning; XLM-RoBERTa; online reviews; AI-generated text; reproducibility.
\end{abstract}

\section{Introduction}

Review texts are useful for monitoring product and service quality, but automated analysis faces at least two distinct questions. The first is what evaluation the author expresses. The second is whether the text belongs to the provenance category for which a system was trained. Generative models make the second question more salient: a plausible review may have been written by a person or generated by an AI system. These properties should not be conflated. A negative rating does not make a review suspicious, and a high classifier score does not establish deception.

In this study, provenance is operationalized as the binary label supplied by MAiDE-up: a human-authored hotel review versus an AI-generated hotel review. MAiDE-up contains real and AI-generated hotel reviews in ten languages; its authors explicitly study how detection performance varies with language, location, and sentiment~\cite{ignat2025maide}. Accordingly, this paper avoids broader claims about ``fake-review detection.''

Multi-task learning assumes that related tasks can use a shared representation and thereby provide one another with an inductive bias~\cite{caruana1997multitask}. For a multilingual corpus, XLM-RoBERTa (XLM-R) is a natural shared encoder because it was pretrained across 100 languages~\cite{conneau2020xlmr}. A shared architecture, however, does not guarantee transfer: an auxiliary task may fail to help, and a single-task comparison is required to separate the effect of shared representations from the effect of using a stronger Transformer model.

The aim of this study is to compare three classifier families on one documented corpus and to constrain the conclusions to the design of that experiment. The research questions are:
\begin{enumerate}[label=\textbf{RQ\arabic*.}, leftmargin=1.5cm]
  \item Does multi-task XLM-R improve Macro-F1 over TF--IDF with logistic regression on either task under the fixed test protocol?
  \item Does multi-task XLM-R differ from a single-task XLM-R trained for the same task?
  \item Which data and experimental properties limit transfer of the results beyond this protocol?
\end{enumerate}

The contribution is not a claim of a universal detector. It is a documented experiment with explicit data provenance, split rules, training configuration, multi-seed results, paired statistical comparisons, and validity limitations.

\section{Related Work}

The idea of training several tasks with shared parameters originates in the formulation of inductive transfer through a common representation~\cite{caruana1997multitask}. In an applied setting, one encoder receives gradients from several classification heads. Whether this is beneficial is an empirical question that depends on task relatedness, label volume, class balance, and the training procedure.

XLM-R was designed for cross-lingual transfer and pretrained as a multilingual masked-language model on data from 100 languages~\cite{conneau2020xlmr}. This motivates its use in a corpus where 23,734 of 30,000 records are in Russian and the remainder covers nine additional languages. The motivation is not evidence of cross-lingual robustness: language slices are not independent external tests in the present study.

The provenance task also depends on the source of the labels. MAiDE-up was created to distinguish genuine from AI-generated hotel reviews and contains 10,000 examples of each type across ten languages~\cite{ignat2025maide}. This study uses a quota-controlled subset of 7,000 texts. The provenance result therefore concerns the synthetic-text setting and domain defined by this resource, not all forms of human deception.

Finally, comparing models on one test set requires more than reporting a mean. Guidance on statistical testing in NLP emphasizes that the choice of test depends on the experimental design and the fixed evaluation examples~\cite{dror2018hitchhiker}. We therefore report variation across random seeds, confidence intervals for paired test-example differences, and an approximate randomization test.

\section{Data and Method}

\subsection{Corpus and Labels}

The input file \repo{data/processed/joint\_reviews.article30k.jsonl} contains 30,000 normalized records. Its SHA-256 digest is \repo{a296cf7bb8cb0dbb3eebd8a0cc9c478422e9f5959836bc4d20915f9f81c9c5fc}. Source roles are summarized in Table~\ref{tab:data}. RuReviews, the Perekrestok card, and the Wildberries card were used as open sources of consumer reviews~\cite{rureviews,perekrestok,wildberries}; MAiDE-up supplied the binary provenance labels~\cite{ignat2025maide}.

\begin{table}[ht]
\centering
\caption{Corpus composition and label provenance.}
\label{tab:data}
\small
\begin{tabularx}{\linewidth}{@{}l r Y Y@{}}
\toprule
Source & $n$ & Sentiment & Provenance \\
\midrule
MAiDE-up & 7,000 & Source label; negative/positive & 3,500 human (authentic) and 3,500 AI-generated (fake) \\
Perekrestok & 10,000 & Rating heuristic: 1--2 negative, 3 neutral, 4--5 positive & No label \\
RuReviews & 3,000 & Native three-class dataset label & No label \\
Wildberries & 10,000 & Rating heuristic: 1--2 negative, 3 neutral, 4--5 positive & No label \\
\bottomrule
\end{tabularx}
\end{table}

All 30,000 records have a sentiment label: 23,734 texts are in Russian and 6,266 are in Turkish, French, Korean, English, Chinese, Spanish, Italian, German, or Romanian. The domain composition is 23,000 e-commerce and 7,000 hospitality records. Provenance labels are present only for the 7,000 MAiDE-up records. The multi-task loss therefore uses each task only where its label is available: e-commerce examples contribute to sentiment training but do not add observations to the provenance loss.

Two components of the sentiment labels are derived from ratings rather than independent manual annotation. This is a deliberate compromise between material volume and target-label precision. The resulting task measures agreement with the project's three-class rating mapping, not a pure linguistic notion of polarity in every possible use of that term.

\subsection{Splitting and Leakage Control}

With split seed 42, the corpus was divided into 23,975 training, 3,013 validation, and 3,012 test records. Before splitting, texts were case-folded and whitespace-normalized, then grouped by normalized text. Groups were stratified by source and the pair of available labels. Consequently, the same normalized review body cannot appear in different partitions.

The audit found 521 exact and 564 normalized duplicate rows in the merged corpus, but zero cross-partition overlap by record identifier, exact text, or normalized text for every train/validation/test pair. Duplicates remain a property of the source material, but the split did not produce observable text leakage. The test sentiment distribution is 710 negative, 204 neutral, and 2,098 positive; the provenance test is balanced at 350 authentic and 350 fake examples. Macro-F1 is therefore the primary metric, with accuracy reported as a secondary metric.

\subsection{Models}

The classical classifier uses TF--IDF uni- and bigrams (up to 50,000 features) with a separate class-balanced LogisticRegression model for each available task. It provides a strong shallow reference on the same training texts.

The single-task and multi-task models use \repo{FacebookAI/xlm-roberta-base}. In the multi-task system, the final hidden state of the first token (CLS) is passed to two independent MLP heads: one for three sentiment classes and one for two provenance classes. Each head has the form $h\rightarrow h\rightarrow k$, GELU nonlinearity, and dropout 0.1. The loss is the sum of weighted cross-entropies for available tasks; both task weights are 0.5. Single-task systems use the same encoder and MLP head but optimize only their own label.

The configuration targets a local Apple Silicon environment: maximum sequence length 192, batch size 2, four-step gradient accumulation, gradient checkpointing, AdamW with learning rate $2\cdot10^{-5}$, weight decay 0.01, 10\% linear warmup, a five-epoch maximum, and early stopping with patience 2. Multi-task training uses a source-balanced sampler and class-weighted cross-entropy; the single-task training loop does not use the source-balanced sampler. Thus, the multi-task versus single-task comparison evaluates two configured training systems, not one isolated architectural switch.

Training ran on a Mac mini M2 Pro with 16~GB unified memory through the MPS backend. The recorded environment was Python 3.12.12 on macOS 26.2 arm64, with MPS as the effective device and training seeds $\{11,21,42,84,126\}$. Checkpoints were selected by validation Macro-F1; the test set was reserved for final evaluation.

\subsection{Comparison Protocol}

For every seed, the multi-task model and both single-task Transformer variants were trained. The classical baseline is deterministic under the fixed configuration, so its standard deviation over repeated report rows is zero; this should not be read as a general measure of method stability. For each Transformer, we report the mean and standard deviation over five seeds and a 95\% t-interval for the mean.

Paired comparisons used the same test examples: 2,000 bootstrap samples for a confidence interval of the seed-averaged difference and an approximate randomization test with 2,000 permutations. The eight checks (two tasks, two metrics, and two comparison families) form a related test family. We therefore interpret direction, interval, and the conservative wording together rather than treating one p-value as decisive.

\section{Results}

\subsection{Test-set Performance}

Table~\ref{tab:mainresults} gives the final metrics. Multi-task XLM-R exceeds TF--IDF with logistic regression on Macro-F1 for both tasks, by 0.0548 for sentiment and 0.0528 for provenance. Single-task Transformer systems have slightly higher means, but that observation alone does not establish a difference: the seed variation is comparable to the mean gap.

\begin{table}[ht]
\centering
\caption{Test metrics; Transformer values are mean $\pm$ standard deviation over five seeds.}
\label{tab:mainresults}
\small
\begin{tabularx}{\linewidth}{@{}Y Y r r@{}}
\toprule
Family & Task & Macro-F1 & Accuracy \\
\midrule
TF--IDF + LR & Sentiment & 0.7161 & 0.8320 \\
Single-task XLM-R & Sentiment & $0.7721\pm0.0087$ & $0.8958\pm0.0038$ \\
Multi-task XLM-R & Sentiment & $0.7709\pm0.0087$ & $0.8950\pm0.0047$ \\
\addlinespace
TF--IDF + LR & Provenance & 0.8998 & 0.9000 \\
Single-task XLM-R & Provenance & $0.9571\pm0.0088$ & $0.9571\pm0.0087$ \\
Multi-task XLM-R & Provenance & $0.9526\pm0.0094$ & $0.9526\pm0.0094$ \\
\bottomrule
\end{tabularx}
\end{table}

For sentiment, the multi-task Macro-F1 t-interval is $[0.7600;0.7817]$, while the single-task interval is $[0.7613;0.7830]$. For provenance, the corresponding intervals are $[0.9409;0.9642]$ and $[0.9463;0.9680]$. These intervals describe training-seed variation and uncertainty under one split; they do not quantify uncertainty from selecting new sources or new text samples.

\subsection{Paired Comparisons}

The main Macro-F1 comparisons appear in Table~\ref{tab:comparisons}. Against the classical baseline, the 95\% bootstrap intervals are positive and approximate randomization p-values are 0.0005, the smallest resolvable value with 2,000 permutations. Against the single-task Transformer, both intervals cross zero and the p-values do not support rejecting no difference. In particular, the multi-task variant did not demonstrate positive transfer; its mean is slightly lower on both tasks.

\begin{table}[ht]
\centering
\caption{Paired Macro-F1 comparisons for the multi-task model. $\Delta=A-B$; intervals use the shared test set.}
\label{tab:comparisons}
\small
\begin{tabularx}{\linewidth}{@{}Y Y r Y r@{}}
\toprule
Task & Comparison ($A$ vs. $B$) & $\Delta$ & 95\% CI & $p_{AR}$ \\
\midrule
Sentiment & multi-task -- single-task & $-0.0013$ & $[-0.0105;0.0081]$ & 0.8016 \\
Sentiment & multi-task -- TF--IDF+LR & $+0.0548$ & $[0.0337;0.0732]$ & 0.0005 \\
Provenance & multi-task -- single-task & $-0.0046$ & $[-0.0109;0.0020]$ & 0.1829 \\
Provenance & multi-task -- TF--IDF+LR & $+0.0528$ & $[0.0305;0.0749]$ & 0.0005 \\
\bottomrule
\end{tabularx}
\end{table}

The answer to RQ1 is positive within the fixed test protocol: multi-task XLM-R has higher Macro-F1 than the selected classical baseline on both tasks. The answer to RQ2 is different: the evidence does not establish a difference between multi-task and single-task XLM-R. Practically, one shared encoder can serve two heads without a clear loss in average quality relative to two single-task Transformer runs, but it does not demonstrate superiority over them.

\section{Discussion and Limitations}

The pattern is asymmetric by comparison level rather than by task. Transformer variants consistently exceed the selected TF--IDF baseline, whereas the difference between shared and task-specific architectures is small relative to the uncertainty estimate. This pattern is compatible with several explanations: the tasks may use partially shared cues, but the volume and type of provenance labels may be insufficient for a measurable multi-task gain. The experiment cannot distinguish these explanations, so no causal claim about the transfer mechanism is made.

The study has six important limitations:
\begin{enumerate}[leftmargin=0.65cm]
  \item \textbf{Limited provenance validity.} All 700 provenance test labels come from MAiDE-up and the hotel domain, contrasting human and AI-generated reviews. This is not an external e-commerce test and not a test of human-written deceptive reviews.
  \item \textbf{Heterogeneous sentiment targets.} Perekrestok and Wildberries sentiment is derived from ratings. Domain differences may therefore reflect the label-generation procedure as well as language and writing style.
  \item \textbf{One fixed split.} Five seeds measure sensitivity to initialization but do not replace repeated splits, leave-one-source-out evaluation, or an external test set. A stronger follow-up should hold out an entire source, domain, or language.
  \item \textbf{Duplicates and sampling.} Grouped splitting removed observable cross-partition leakage, but 564 normalized duplicate rows remain in the merged corpus. Sources were also included by fixed quotas, so the aggregate metric is not an estimate for a random review population.
  \item \textbf{Non-isolated ablation.} The multi-task loop uses a source-balanced sampler while the single-task loop does not. A pure architectural ablation requires repeating the comparison with identical sampling schedules.
  \item \textbf{No calibration or deployment validation.} Softmax probabilities were not calibrated, a human-review threshold was not selected, and the checkpoints were not automatically promoted as the application's production model. The metrics therefore do not define a moderation policy.
\end{enumerate}

These limitations determine the next research steps. First, independent provenance labels should be collected for another domain and for human-deception data. Second, source-held-out and language-held-out protocols should be repeated with identical sampling for the multi-task and single-task systems. Third, an applied system needs calibration, a human-in-the-loop interface, and error analysis rather than an averaged F1 score alone.

\section{Conclusion}

This paper implemented and evaluated a reproducible shared XLM-RoBERTa encoder for sentiment and text provenance on a partially labelled mixed corpus of 30,000 reviews. Multi-task training increased test Macro-F1 over TF--IDF with logistic regression from 0.7161 to 0.7709 for sentiment and from 0.8998 to 0.9526 for provenance. Single-task Transformer systems nevertheless produced practically comparable and slightly higher means, and paired intervals did not confirm a difference.

The result supports a shared model as a compact research artifact, but it does not support a claim of proven superiority over a single-task Transformer or universal fake-review detection. The most important next step is an external, source-held-out and domain-held-out evaluation with independent provenance labels.

\section*{Data Availability and Reproducibility}

Scripts, configuration, and aggregated reports are available in the local project repository: \repo{configs/model.article30k.mac\_m2\_16gb.yaml}, \repo{reports/multiseed/article30k-m2pro-retry1/}, \repo{reports/multiseed/article30k-m2pro-retry1\_statistics/}, and \repo{reports/audit/joint\_reviews.article30k.audit.json}. The processed-corpus digest is given in Section~3.1. Redistribution of source records is governed by the original dataset terms; the Wildberries component is explicitly CC BY-NC-SA 4.0. Exact reproduction should use the recorded configuration snapshot and the cited dataset cards rather than silently substituting newer versions.

\section*{Ethics Statement}

No participants were recruited and no new personal data were collected. The analysis uses textual records obtained from open research sources; the manuscript does not reproduce raw reviews, author names, or direct user identifiers. Provenance predictions should be treated as research signals, not as standalone evidence of author misconduct.

\section*{CRediT Author Contributions}

\textbf{Replace this template with actual names and roles before submission.} For a single-author version, the expected statement is: \textit{[Author]: conceptualization, methodology, software, data curation, formal analysis, validation, visualization, writing---original draft, and writing---review and editing.}

\section*{Funding and Conflict of Interest}

Funding and conflict-of-interest information was not supplied by the author. Before submission, replace this text with an explicit declaration, for example ``No external funding was received'' and ``The author declares no conflict of interest,'' or provide the applicable details.

\section*{AI Use Disclosure}

Generative AI was used as an auxiliary tool for manuscript structure and language drafting. Numerical results, experimental configuration, data claims, and references were checked against the repository artifacts and primary sources. The author remains responsible for the final manuscript, interpretation, and compliance with the submission venue's policies.

\bibliographystyle{unsrtnat}
\bibliography{article30k_m2pro_ru}

\end{document}
